\documentclass[12pt]{article}

\usepackage{amsmath, amssymb, amsthm}
\usepackage{geometry}
\usepackage{setspace}

\geometry{margin=1in}
\onehalfspacing

\title{CSE 400: Fundamentals of Probability in Computing \\ \large Tutorial--1 Lecture Scribe}
\author{Name: Masoom Choksi \\ AUID: AU2440263}
\date{February 13, 2026}

\begin{document}
\maketitle

\section*{Question 1}

\textbf{Problem Statement.}  
There are 20 distinct dishes to be divided into 4 groups of 5 dishes each. The order of the groups does not matter.

\subsection*{(a) Total number of ways}

\textbf{Assumptions.}
\begin{itemize}
\item All dishes are distinct.
\item Groups are unordered.
\item Each group has exactly 5 dishes.
\end{itemize}

\textbf{Derivation.}

The total number of ways to partition 20 distinct objects into 4 groups of equal size 5 is

\[
\frac{20!}{(5!)^4 \cdot 4!}
\]

The denominator accounts for:
\begin{itemize}
\item $5!$ permutations within each group,
\item $4!$ permutations of the groups themselves.
\end{itemize}

\subsection*{(b) Probability all platters have the same cuisine}

\textbf{Given.}
\begin{itemize}
\item 5 dishes from each of 4 cuisines.
\end{itemize}

\textbf{Favorable outcomes.}  
Each platter consists of dishes from exactly one cuisine. Since cuisines are fixed, there is only

\[
1
\]

favorable arrangement.

\textbf{Total outcomes.}
\[
\frac{20!}{(5!)^4 \cdot 4!}
\]

\textbf{Probability.}
\[
P = \frac{1}{\frac{20!}{(5!)^4 \cdot 4!}} = \frac{(5!)^4 \cdot 4!}{20!}
\]

\section*{Question 2}

\textbf{Problem Statement.}  
A passenger arrives uniformly between 7:00 and 7:30. Buses arrive every 15 minutes.

\textbf{Assumption.}
Arrival time is uniformly distributed over a 30-minute interval.

\subsection*{(a) Waiting time less than 5 minutes}

Total interval length:
\[
30 \text{ minutes}
\]

Favorable intervals:
\[
2 \times 5 = 10 \text{ minutes}
\]

\textbf{Probability.}
\[
P = \frac{10}{30} = \frac{1}{3}
\]

\subsection*{(b) Waiting time more than 10 minutes}

Favorable intervals:
\[
2 \times 5 = 10 \text{ minutes}
\]

\textbf{Probability.}
\[
P = \frac{10}{30} = \frac{1}{3}
\]

\section*{Question 3}

\textbf{Given.}
\[
P(L) = 0.15, \quad P(E) = 0.25, \quad P(L \cap E) = 0.08
\]

\textbf{Required.}
\[
P(L^c \mid E^c)
\]

\textbf{Derivation.}
\[
P(E^c) = 1 - 0.25 = 0.75
\]

\[
P(L^c \cap E^c) = 1 - P(L \cup E)
\]

\[
P(L \cup E) = 0.15 + 0.25 - 0.08 = 0.32
\]

\[
P(L^c \cap E^c) = 1 - 0.32 = 0.68
\]

\textbf{Conditional Probability.}
\[
P(L^c \mid E^c) = \frac{0.68}{0.75} = \frac{68}{75}
\]

\section*{Question 4}

\textbf{Given.}  
Number of typographical errors follows a Poisson distribution with
\[
\lambda = 0.2
\]

\subsection*{(a) Probability of 0 errors}

\[
P(X=0) = e^{-0.2}
\]

\subsection*{(b) Probability of 2 or more errors}

\[
P(X \ge 2) = 1 - P(0) - P(1)
\]

\[
P(1) = 0.2 e^{-0.2}
\]

\[
P(X \ge 2) = 1 - e^{-0.2} - 0.2 e^{-0.2}
\]

\section*{Question 5}

\textbf{Given.}
\[
\lambda = 3.5
\]

\subsection*{(a) At least 2 accidents}

\[
P(X \ge 2) = 1 - P(0) - P(1)
\]

\[
P(X \ge 2) = 1 - e^{-3.5} - 3.5 e^{-3.5}
\]

\subsection*{(b) At most 1 accident}

\[
P(X \le 1) = e^{-3.5} + 3.5 e^{-3.5}
\]

\section*{Question 6}

\textbf{Given.}
\[
X \sim \text{Poisson}(3)
\]

\subsection*{(a) $P(X \ge 3)$}

\[
P(X \ge 3) = 1 - (P(0) + P(1) + P(2))
\]

\[
= 1 - \left(e^{-3} + 3e^{-3} + \frac{9}{2}e^{-3}\right)
\]

\subsection*{(b) $P(X \ge 3 \mid X \ge 1)$}

\[
P(X \ge 3 \mid X \ge 1) = \frac{P(X \ge 3)}{1 - P(0)}
\]

\section*{Question 7}

\textbf{Given PDF.}
\[
f_X(x) = \frac{1}{\sqrt{8\pi}} \exp\left(-\frac{(x+3)^2}{8}\right)
\]

\textbf{Standardization.}
\[
Z = \frac{X+3}{2}
\]

\subsection*{Results (Q-function form)}

\[
P(X \le 0) = Q\left(\frac{3}{2}\right)
\]

\[
P(X > 4) = Q\left(\frac{7}{2}\right)
\]

\[
P(|X+3|<2) = 1 - 2Q(1)
\]

\[
P(|X-2|>1) = Q\left(\frac{1}{2}\right) + Q\left(\frac{3}{2}\right)
\]

\section*{Question 8}

\textbf{Given.}
\begin{itemize}
\item 12 jurors
\item At least 8 must vote correctly
\item Each juror votes correctly with probability $\theta$
\end{itemize}

\textbf{Probability of correct decision.}

\[
P = \sum_{k=8}^{12} \binom{12}{k} \theta^k (1-\theta)^{12-k}
\]

\section*{Question 9}

\textbf{PMF.}
\[
p(i) = \frac{c\lambda^i}{i!}
\]

\textbf{Normalization.}
\[
\sum_{i=0}^{\infty} \frac{c\lambda^i}{i!} = 1
\Rightarrow c = e^{-\lambda}
\]

\subsection*{(a) $P(X=0)$}

\[
P(X=0) = e^{-\lambda}
\]

\subsection*{(b) $P(X>2)$}

\[
P(X>2) = 1 - e^{-\lambda}\left(1 + \lambda + \frac{\lambda^2}{2}\right)
\]

\section*{Question 10}

\textbf{Given.}
\begin{itemize}
\item $n$ components
\item Each functions with probability $p$
\item System works if at least half function
\end{itemize}

\subsection*{(a) Comparison: 5 vs 3 components}

\[
P_5 = \sum_{k=3}^{5} \binom{5}{k} p^k (1-p)^{5-k}
\]

\[
P_3 = \sum_{k=2}^{3} \binom{3}{k} p^k (1-p)^{3-k}
\]

System with 5 components is better when:
\[
P_5 > P_3
\]

\subsection*{(b) General case}

A $(2k+1)$-component system is better than $(2k-1)$ when:
\[
p > \frac{1}{2}
\]

\end{document}